\documentclass[11pt]{article}

\usepackage[utf8]{inputenc}
\usepackage[T1,T2A]{fontenc}
\usepackage{lmodern}
\usepackage[english,ukrainian]{babel}
\selectlanguage{english}
\usepackage[margin=1in]{geometry}
\usepackage{amsmath,amssymb}
\usepackage{booktabs}
\usepackage{multirow}
\usepackage{array}
\usepackage{tabularx}
\usepackage{graphicx}
\usepackage{float}
\usepackage{tikz}
\usepackage[dvipsnames]{xcolor}
\usepackage{enumitem}
\usepackage{caption}
\usepackage{subcaption}
\usepackage{microtype}
\usepackage{url}
\usepackage[numbers,sort&compress]{natbib}
\usepackage{hyperref}
\usepackage{pgfplots}
\usepackage{pgfplotstable}
\pgfplotsset{compat=1.18}
\usepgfplotslibrary{colorbrewer}
\usetikzlibrary{arrows.meta,positioning,calc,shapes.geometric}

\hypersetup{
  colorlinks=true,
  linkcolor=RoyalBlue,
  citecolor=ForestGreen,
  urlcolor=BrickRed
}

\newcolumntype{R}[1]{>{\raggedleft\arraybackslash}p{#1}}

\title{DefectRadar: Automated Detection of Definitional Defects\\in Legislation via Morpho-Semantic NLP Pipeline\\over 24{,}000 Ukrainian Laws}

\author{
  Volodymyr Ovcharov$^{1}$ \quad Ihor Kyrychenko$^{2}$ \\[6pt]
  $^{1}$LEX AI LLC, Kyiv, Ukraine \\
  $^{2}$Independent Legal Researcher, Kyiv, Ukraine \\[4pt]
  \texttt{vladimir@legal.org.ua}
}

\date{}

\begin{document}
\renewcommand{\abstractname}{Abstract}
\maketitle

\begin{abstract}
Legal definitions in legislation are assumed to satisfy basic logical requirements -- yet we find that \textbf{52\% of all statutory definitions} in Ukrainian law contain at least one formally detectable logical defect.
A definition must not define a term through itself (\emph{circulus in definiendo}) nor rely on terms that lack their own definitions (\emph{ignotum per ignotum}); violations produce normatively empty provisions that undermine legal certainty.
Prior NLP work has addressed definition extraction from legislative text~\citep{hosabettu2025statutory}, stylistic quality checking of legislative drafts~\citep{hofler2012style}, and retrieval-augmented definition generation~\citep{chouhan2024lexdrafter}, but no system has attempted to detect \emph{logical} defects within definitions themselves.
We present \textsc{DefectRadar}, the first pipeline for automated detection of \emph{circulus in definiendo} and \emph{ignotum per ignotum} in legislation.
The pipeline extracts 5{,}799 definitions from 879 laws within a corpus of 44{,}021 active Ukrainian legislative acts, applies morpho-semantic circularity detection (Ukrainian lemmatization via pymorphy2-uk combined with discriminating-adjective analysis), and constructs a cross-legislative definition graph to identify undefined-term dependencies across the entire corpus.
DefectRadar identifies 1{,}880 circular definitions (32.4\%) and 1{,}885 cross-reference gaps (32.5\%), with the highest defect concentration in post-2014 EU harmonization laws -- electronic communications, media, payment services, pharmaceuticals, and insurance -- consistent with a ``decorative borrowing'' pattern in \emph{acquis} transposition.
As a diagnostic complement to generative systems like LexDrafter, DefectRadar enables a detect--diagnose--repair workflow for legislative definition quality.
We release the pipeline code and audit results at \url{https://huggingface.co/datasets/overthelex/ua-defectradar}.
\end{abstract}

\noindent\textbf{Keywords:} legal NLP, legislative quality assurance, definitional defects, \emph{circulus in definiendo}, \emph{ignotum per ignotum}, Ukrainian law, morphological analysis, cross-legislative definition graph

% ============================================================
\section{Introduction}
\label{sec:intro}
% ============================================================

Every modern legal system relies on definitions.
When a legislature introduces a term -- ``creative industries,'' ``beneficial owner,'' ``critical infrastructure'' -- it is expected to provide a definition that unambiguously delineates the concept's scope.
The rules governing such definitions have been stable since Aristotle's \emph{Topics}: a definition must specify the proximate genus and differentiating characteristics (\emph{definitio per genus proximum et differentiam specificam}), must not contain the term being defined (\emph{circulus in definiendo}), and must not explain the unknown through equally unknown concepts (\emph{ignotum per ignotum})~\citep{hart2012concept,dzeiko2007legislative}.

Violations of these rules are not merely stylistic.
A tautological definition -- one that defines ``creative industries'' as economic activities pursued through ``creative expression'' -- produces a normatively empty provision: it imposes no constraint on the concept's application and delegates unlimited discretion to the executive branch.
\citet{kyrychenko2026creative} demonstrated this concretely for Ukrainian law, showing that a single \emph{circulus in definiendo} in the Law on Culture enabled hundreds of millions of hryvnias in discretionary grants to flow through a legally undefined category over six years.

The natural question is whether this case is isolated or systemic.
Ukraine maintains over 24{,}000 active legislative acts, each potentially containing dozens of definitions.
Manual auditing of this corpus is infeasible -- even expert legal analysis of a single defective definition, as performed by \citet{kyrychenko2026creative}, requires months of doctrinal work.
The same scaling problem affects every jurisdiction with a substantial legislative corpus: the EU \emph{acquis}, German federal law, or the US Code.

Prior computational work on legislative quality has addressed stylistic checking~\citep{hofler2012style}, definition extraction~\citep{hosabettu2025statutory}, and definition generation~\citep{chouhan2024lexdrafter}, but no system has attempted to detect \emph{logical} defects within definitions themselves.
We address this gap with \textsc{DefectRadar}, the first automated pipeline for detecting definitional defects in legislation.
Our contributions are:

\begin{enumerate}[nosep,leftmargin=*]
  \item A \textbf{definition extraction module} that exploits the structural regularity of Ukrainian legislative text to extract 5{,}799 legal definitions from 879 laws within a corpus of 44{,}021 active legislative acts.
  \item A \textbf{circulus in definiendo detector} that combines Ukrainian morphological analysis (lemmatization via pymorphy2-uk) with discriminating-adjective extraction to detect tautological definitions where qualifying terms from the \emph{definiendum} reappear in the \emph{differentia specifica}.
  \item A \textbf{cross-legislative ignotum detector} that constructs a definition graph across the entire corpus and identifies terms used in \emph{definientia} that are defined in other laws but not in the current one -- revealing cross-reference gaps in legislative drafting.
  \item A \textbf{full-corpus audit} of active Ukrainian legislation revealing 1{,}880 \emph{circulus} defects (32.4\%) and 1{,}885 cross-reference \emph{ignotum} gaps (32.5\%), with 749 definitions exhibiting both defects simultaneously (12.9\%).
  \item A \textbf{domain analysis} showing that EU harmonization laws (electronic communications, media, payment services, pharmaceuticals, insurance) concentrate the highest defect densities -- consistent with the ``decorative borrowing'' pattern identified by \citet{kyrychenko2026creative}.
\end{enumerate}

The pipeline is language-aware but architecture-agnostic: we describe adaptation requirements for other morphologically rich languages (Polish, Czech, Lithuanian) whose legislative corpora share structural conventions with Ukrainian law.

% ============================================================
\section{Background: Definitional Defects in Legislation}
\label{sec:background}
% ============================================================

Before describing our NLP approach, we establish the legal-theoretical framework that defines what constitutes a definitional defect.
This framework draws on the doctrinal analysis of \citet{kyrychenko2026creative}, which we formalize for computational purposes.

\subsection{Anatomy of a Legal Definition}
\label{sec:anatomy}

A legal definition in Ukrainian legislation follows the canonical form inherited from Roman legal tradition:

\begin{quote}
\emph{definiendum} -- \emph{genus proximum}, \emph{differentia specifica}.
\end{quote}

\noindent For example, Article~1 of the Law of Ukraine ``On Culture'' defines:

\begin{quote}
\foreignlanguage{ukrainian}{креативні індустрії -- види економічної діяльності, метою яких є створення доданої вартості і робочих місць через культурне (мистецьке) та/або креативне вираження.}
\end{quote}

\noindent Here the \emph{definiendum} is \foreignlanguage{ukrainian}{``креативні індустрії''} (creative industries), the \emph{genus proximum} is \foreignlanguage{ukrainian}{``види економічної діяльності''} (types of economic activity), and the \emph{differentia specifica} is the remainder of the definition.

\subsection{Circulus in Definiendo}
\label{sec:circulus-def}

A definition exhibits \emph{circulus in definiendo} when the \emph{definiens} contains the \emph{definiendum} or its morphological derivative.
In the example above, the adjective \foreignlanguage{ukrainian}{``креативний''} (creative) in the \emph{definiendum} reappears as \foreignlanguage{ukrainian}{``креативне вираження''} (creative expression) in the \emph{definiens}.
We formalize this as:

\begin{equation}
\label{eq:circulus}
\text{Circulus}(d) = \begin{cases}
1 & \text{if } \exists\, t \in \text{lemmas}(D_{\text{dum}}) : t \in \text{lemmas}(D_{\text{ens}}) \setminus \text{lemmas}(G) \\
0 & \text{otherwise}
\end{cases}
\end{equation}

\noindent where $D_{\text{dum}}$ is the \emph{definiendum}, $D_{\text{ens}}$ is the \emph{definiens}, $G$ is the \emph{genus proximum}, and $\text{lemmas}(\cdot)$ returns the set of content-word lemmas (excluding stop words and function words).
The exclusion of $G$ is critical: many definitions legitimately repeat domain-generic terms (e.g., \foreignlanguage{ukrainian}{``діяльність''}, activity) that appear in both the term being defined and the genus; these are not tautological.

This lexical test captures \emph{explicit} circularity.
We augment it with a semantic test using sentence embeddings to detect \emph{implicit} circularity -- cases where the \emph{definiens} is semantically equivalent to a paraphrase of the \emph{definiendum} without sharing surface-level morphemes.

\subsection{Ignotum per Ignotum}
\label{sec:ignotum-def}

A definition exhibits \emph{ignotum per ignotum} when the \emph{definiens} relies on terms that are themselves undefined in the legislative corpus.
This is distinct from \emph{circulus}: the definition may be non-circular but still empty because it explains the unknown through the equally unknown.

Formally, let $\mathcal{D}$ be the set of all defined terms across the legislative corpus.
A definition $d$ with \emph{definiens} $D_{\text{ens}}$ exhibits \emph{ignotum per ignotum} if:

\begin{equation}
\label{eq:ignotum}
\exists\, t \in \text{terms}(D_{\text{ens}}) : t \notin \mathcal{D} \land t \in \mathcal{T}_{\text{legal}}
\end{equation}

\noindent where $\text{terms}(\cdot)$ extracts noun phrases and $\mathcal{T}_{\text{legal}}$ is the set of domain-specific legal terms (as opposed to common-language words).
The restriction to $\mathcal{T}_{\text{legal}}$ prevents false positives from everyday words like \foreignlanguage{ukrainian}{``особа''} (person) or \foreignlanguage{ukrainian}{``майно''} (property) that have well-established legal meanings even without explicit statutory definitions.

% ============================================================
\section{Related Work}
\label{sec:related}
% ============================================================

\subsection{Legal NLP}

The past five years have seen rapid progress in NLP for legal text.
Comprehensive surveys by \citet{katz2023natural} and \citet{ariai2024survey} catalog tasks including document classification, named entity recognition, judgment prediction, and argument mining.
Benchmarks such as LexGLUE~\citep{chalkidis2022lexglue} and LEXTREME~\citep{niklaus2023lextreme} standardize evaluation across jurisdictions and languages.
The MultiLegalPile corpus~\citep{niklaus2023multilegal} provides 689\,GB of multilingual legal text for pretraining.
LEMUR~\citep{ahmadi2026lemur} introduces multilingual legal embeddings optimized for retrieval.

However, all of these resources focus on \emph{understanding} legal text as written.
None address the orthogonal problem of detecting \emph{defects} in how legal text is written -- i.e., legislative quality assurance.

\subsection{Legal Definition Extraction}

Automatic extraction of legal definitions is a prerequisite for any downstream quality analysis.
\citet{hosabettu2025statutory} present a transformer-based pipeline for extracting statutory definitions from the U.S.\ Code, achieving 96.8\% precision and 98.9\% recall using Legal-BERT fine-tuned on the XML structure of federal statutes.
Their multi-stage architecture -- document structure analysis followed by neural classification -- parallels our Stage~1 extractor.
However, the U.S.\ Code provides explicit XML markup (\texttt{<term>} tags, section-level definition scope annotations) that Ukrainian legislation lacks entirely: Verkhovna Rada publishes legislative text as unstructured HTML without semantic tagging.
This structural gap explains our hybrid approach (Section~\ref{sec:extraction}), which combines regex-based exploitation of typographic conventions (the em-dash copula) with a fine-tuned classifier for inline definitions.
Crucially, \citeauthor{hosabettu2025statutory}'s work stops at extraction -- it does not analyze the \emph{quality} of extracted definitions, leaving circulus and ignotum defects undetected.

In a complementary direction, \citet{chouhan2024lexdrafter} introduce LexDrafter, a retrieval-augmented generation framework that \emph{drafts} definitions for new EU legislative documents by retrieving and synthesizing existing definitions from the \emph{acquis}.
LexDrafter constructs a cross-legislative definition knowledge base from EU energy-domain regulations and uses RAG to generate harmonized definitions on demand.
This is architecturally close to our Stage~3 definition graph: both systems build inter-legislative term dependency structures.
The key difference is directional -- LexDrafter is \emph{generative} (producing new definitions), while our Stage~3 is \emph{diagnostic} (identifying gaps in existing ones).
The two approaches are complementary: DefectRadar could identify cross-reference \emph{ignotum} gaps, and a LexDrafter-style system could then propose definitions to fill them.

\subsection{Norm Conflict and Fallacy Detection}

\citet{aires2019classifying} use learned semantic representations to classify conflicts between normative statements in Brazilian legislation.
\citet{holzenberger2020statutory} construct a statutory reasoning dataset for US tax law, though their focus is on entailment rather than structural defects.
\citet{sadowski2025verifiable} propose a multi-agent framework for formal legal reasoning with symbolic verification.
Contract-level NLI~\citep{koreeda2021contractnli} detects contradictions between contract clauses and hypotheses.
\citet{nakpih2020fallacies} develop a Prolog-based system for automated discovery of logical fallacies (specifically \emph{non sequitur}) in legal argumentation, checking validity and soundness of judicial claims.

These approaches share our goal of automated quality assurance for legal text, but differ in scope.
Norm conflict detection addresses inconsistencies \emph{between} norms; fallacy detection targets reasoning errors in \emph{judicial argumentation}.
Our work addresses defects \emph{within} individual legislative norms -- specifically, within the logical structure of definitions.
A definition can be internally defective (circular or reliant on undefined terms) while being perfectly consistent with all other norms in the corpus and free of argumentative fallacies.

\subsection{Legislative Drafting Quality}

The legal scholarship on legislative technique is extensive~\citep{hart2012concept,mccormick1991interpreting,kelsen2004pure,dzeiko2007legislative,onischuk2014technique}, but computational approaches are scarce.
The earliest work we are aware of is \citet{hofler2012style}, who develop an automated style checker for Swiss German-language legislative drafts.
Their system detects violations of drafting guidelines -- passive voice, overly long sentences, ambiguous modal verbs, inconsistent terminology -- by adapting error-modeling techniques from controlled-language checkers for technical writing.
Crucially, \citeauthor{hofler2012style} explicitly identify \emph{definitional} rules as targets for automated checking, including ``a term must not be defined by itself'' -- which is precisely \emph{circulus in definiendo}.
However, they classify this as an ``abstract rule'' requiring deep ``linguistic concretisation'' before it can be assessed by machine, and their implemented system checks only stylistic violations, leaving definitional logic unaddressed.
Our Stage~2 provides exactly this concretisation: the morphological lemmatization and genus-exclusion approach translates the abstract rule ``a term must not be defined by itself'' into a computable check over Ukrainian inflectional morphology (Section~\ref{sec:circulus-method}).
DefectRadar thus closes an open problem identified 14 years ago in the legislative NLP literature.

More recently, \citet{braun2024agb} evaluate clause legality in German consumer contracts, and \citet{wardas2025employment} benchmark contract review for employment law.
The CLASE system~\citep{ma2026clase} evaluates stylistic quality of Chinese legal text using a hybrid of linguistic features and LLM-as-judge scoring.
GELATO~\citep{flynn2026gelato} provides legislative NER but does not assess definitional quality.

Table~\ref{tab:related-positioning} summarizes where DefectRadar sits relative to these systems.
To our knowledge, no prior work has attempted automated detection of \emph{circulus in definiendo} or \emph{ignotum per ignotum} in legislation.

\begin{table}[t]
\centering
\caption{Positioning of DefectRadar relative to existing legislative NLP systems.}
\label{tab:related-positioning}
\small
\begin{tabular}{lcccc}
\toprule
\textbf{System} & \textbf{Extract} & \textbf{Style} & \textbf{Logic} & \textbf{Cross-law} \\
\midrule
H{\"o}fler \& Sugisaki '12 & -- & \checkmark & -- & -- \\
Hosabettu \& Shah '25 & \checkmark & -- & -- & -- \\
LexDrafter '24 & \checkmark & -- & -- & \checkmark$^{*}$ \\
CLASE '26 & -- & \checkmark & -- & -- \\
GELATO '26 & \checkmark & -- & -- & -- \\
\textsc{DefectRadar} & \checkmark & -- & \checkmark & \checkmark \\
\bottomrule
\multicolumn{5}{l}{\footnotesize $^{*}$Generative (drafts new definitions), not diagnostic.}
\end{tabular}
\end{table}

\subsection{Ukrainian Legal NLP}

Ukrainian is severely underrepresented in legal NLP.
LEXTREME~\citep{niklaus2023lextreme} covers 24 languages but excludes Ukrainian entirely.
\citet{kiulian2024borsch} fine-tune Gemma and Mistral for Ukrainian but focus on general language understanding rather than legal domain.
\citet{ovcharov2026tokenizer} benchmark seven foundation models on Ukrainian legal text, finding 1.6$\times$ tokenizer fertility spread and counterintuitive few-shot degradation.
\citet{ovcharov2026statute} construct a statute retrieval benchmark from 396M Ukrainian court citations, documenting temporal decay of co-citation structure.
\citet{ovcharov2026temporal} analyze temporal drift in citation patterns across 101M court decisions.

Our work extends Ukrainian legal NLP from analysis and retrieval to \emph{legislative quality assurance} -- a task that requires deep integration of linguistic knowledge (Ukrainian morphology) with legal-theoretical knowledge (definition structure).

% ============================================================
\section{Data}
\label{sec:data}
% ============================================================

\subsection{Legislative Corpus}
\label{sec:corpus}

We construct our corpus from the Verkhovna Rada Open Data portal (\foreignlanguage{ukrainian}{Верховна Рада відкриті дані}, \texttt{data.rada.gov.ua}), which provides machine-readable access to the full text of Ukrainian legislation.
Using the LEX AI platform's RadaLegislationAdapter~\citep{ovcharov2026tokenizer}, we retrieve all legislative acts with status ``active'' (\foreignlanguage{ukrainian}{чинний}) as of April 2026.

The corpus comprises:

\begin{itemize}[nosep,leftmargin=*]
  \item 44{,}021 active legislative acts (1{,}274 laws, 61 codes, 42{,}643 regulations, 43 untyped)
  \item Temporal span: 1991--2026 (Ukrainian independence to present)
  \item All acts include full text (HTML and plain text variants)
  \item Acts range from single-article resolutions to the Civil Code (1{,}308 articles)
\end{itemize}

\subsection{Definition Annotation}
\label{sec:annotation}

For evaluation, we plan a gold-standard annotation of 500 definitions sampled from across the corpus using stratified sampling by legal domain and temporal period.
In the current version, we validate against the known case documented by \citet{kyrychenko2026creative} (the ``creative industries'' definition in the Law on Culture) and perform manual inspection of 100 randomly sampled pipeline outputs.

Each definition in the planned annotation will be labeled along three dimensions:

\begin{enumerate}[nosep,leftmargin=*]
  \item \textbf{Definition boundaries}: start/end offsets of \emph{definiendum}, \emph{genus proximum}, and \emph{differentia specifica}
  \item \textbf{Circulus in definiendo}: binary label (defective / non-defective), with free-text justification identifying the circular element
  \item \textbf{Ignotum per ignotum}: list of terms in the \emph{definiens} that are defined in other laws but absent from the current one (cross-reference gaps)
\end{enumerate}

% ============================================================
\section{Method}
\label{sec:method}
% ============================================================

\textsc{DefectRadar} operates as a three-stage pipeline (Figure~\ref{fig:pipeline}).

\begin{figure}[t]
\centering
\begin{tikzpicture}[
  node distance=1.8cm and 2.2cm,
  stage/.style={draw, rounded corners=3pt, minimum width=3.2cm, minimum height=1.1cm, align=center, font=\small\sffamily},
  input/.style={draw, dashed, rounded corners=2pt, minimum width=2.8cm, minimum height=0.8cm, align=center, font=\footnotesize\sffamily, fill=gray!8},
  output/.style={draw, rounded corners=2pt, minimum width=2.8cm, minimum height=0.8cm, align=center, font=\footnotesize\sffamily, fill=ForestGreen!10},
  arrow/.style={-{Stealth[length=5pt]}, thick},
  label/.style={font=\tiny\sffamily\bfseries, fill=white, inner sep=1pt}
]

% Input
\node[input] (corpus) {Legislative Corpus\\24{,}187 laws};

% Stage 1
\node[stage, right=of corpus, fill=RoyalBlue!12] (s1) {\textbf{Stage 1}\\Definition\\Extraction};

% Stage 2
\node[stage, right=of s1, fill=BrickRed!12] (s2) {\textbf{Stage 2}\\\emph{Circulus}\\Detection};

% Stage 3
\node[stage, right=of s2, fill=Goldenrod!15] (s3) {\textbf{Stage 3}\\\emph{Ignotum}\\Detection};

% Outputs
\node[output, below=1.2cm of s1] (defs) {47{,}312\\definitions};
\node[output, below=1.2cm of s2] (circ) {312 \emph{circulus}\\defects (1.8\%)};
\node[output, below=1.2cm of s3] (igno) {2{,}847 \emph{ignotum}\\defects (16.4\%)};

% Tools
\node[font=\tiny\sffamily, text=gray, below=0.15cm of s1] {regex + XLM-R};
\node[font=\tiny\sffamily, text=gray, below=0.15cm of s2] {pymorphy2 + E5};
\node[font=\tiny\sffamily, text=gray, below=0.15cm of s3] {term graph};

% Arrows
\draw[arrow] (corpus) -- (s1);
\draw[arrow] (s1) -- (s2);
\draw[arrow] (s2) -- (s3);
\draw[arrow] (s1) -- (defs);
\draw[arrow] (s2) -- (circ);
\draw[arrow] (s3) -- (igno);

\end{tikzpicture}

\caption{The \textsc{DefectRadar} pipeline. Stage~1 extracts definitions from legislative text. Stage~2 detects \emph{circulus in definiendo} via morphological and semantic analysis. Stage~3 constructs a cross-legislative definition graph to detect \emph{ignotum per ignotum}.}
\label{fig:pipeline}
\end{figure}

\subsection{Stage 1: Definition Extraction}
\label{sec:extraction}

Ukrainian legislative acts follow rigid structural conventions inherited from Soviet normographic tradition and maintained by the Ministry of Justice's methodological guidelines~\citep{minjust2006guidelines}.
Definitions are concentrated in Article~1 (``Definition of Terms'') of each law, though they also appear inline throughout the text.

Unlike the U.S.\ Code, where \citet{hosabettu2025statutory} leverage explicit XML markup (\texttt{<term>} tags and section-level scope annotations) for definition extraction, Ukrainian legislation is published as unstructured HTML without semantic tagging.
This precludes a purely structure-driven approach and motivates our hybrid extractor:

\begin{enumerate}[nosep,leftmargin=*]
  \item \textbf{Structural rules}: Regex patterns matching the canonical definition format: \texttt{<term> -- <definition>;} or \texttt{<term> -- ce <definition>}. Ukrainian legislative definitions consistently use the em-dash (\foreignlanguage{ukrainian}{«--»}) as the definitional copula -- a typographic convention as reliable as XML markup, but requiring language-specific pattern engineering.
  \item \textbf{Article-1 parser}: A specialized parser for Article~1 blocks that handles multi-sentence definitions, numbered sub-definitions, and cross-references.
  \item \textbf{Inline definition detector}: A classifier (fine-tuned XLM-RoBERTa-base) that identifies inline definitions occurring outside Article~1, trained on 500 manually labeled examples. This component is analogous to \citeauthor{hosabettu2025statutory}'s Legal-BERT classifier, adapted for Ukrainian legislative conventions.
\end{enumerate}

The rule-based components (1--2) handle 94\% of definitions; the neural component (3) catches the remaining inline definitions that lack the canonical em-dash pattern.
This division reflects a deliberate design choice: for well-structured legislative text, rule-based extraction is more interpretable and requires no training data, while the neural fallback ensures coverage of non-canonical patterns.

For each extracted definition, we decompose it into:
\begin{itemize}[nosep,leftmargin=*]
  \item \emph{Definiendum}: the term being defined
  \item \emph{Genus proximum}: the superordinate category (typically the first noun phrase after the copula)
  \item \emph{Differentia specifica}: the remainder of the definition
\end{itemize}

This decomposition goes beyond flat extraction: \citet{hosabettu2025statutory} extract (term, definition, scope) triples, while we further segment the \emph{definiens} into genus and differentia -- a finer granularity required for circulus detection (Stage~2).

\subsection{Stage 2: Circulus in Definiendo Detection}
\label{sec:circulus-method}

The core challenge of automated circulus detection is what \citet{hofler2012style} call \emph{concretisation}: translating the abstract rule ``a term must not be defined by itself'' into a computable check.
The difficulty lies in morphologically rich languages where the same root appears in multiple inflected forms -- Ukrainian has 7 grammatical cases, 3 genders, and productive diminutive/augmentative derivation, so surface-string matching is insufficient.
We implement a two-tier detector:

\paragraph{Tier 1: Morphological overlap.}
We lemmatize both the \emph{definiendum} and the \emph{differentia specifica} using pymorphy2-uk~\citep{korobov2015morphological}, a morphological analyzer for Ukrainian that handles the language's complex inflectional system (7 cases, 3 genders, diminutive/augmentative forms).
We extract content-word lemmas (excluding the lemmas of the \emph{genus proximum} to avoid false positives from legitimate domain terms) and check for lemma overlap:

\begin{equation}
\text{MorphCirculus}(d) = |\text{lemmas}(D_{\text{dum}}) \cap (\text{lemmas}(D_{\text{spec}}) \setminus \text{lemmas}(G))| > 0
\end{equation}

This catches explicit circularity (the ``creative industries'' case), but misses paraphrastic circularity.

\paragraph{Tier 2: Semantic similarity.}
For definitions that pass Tier~1 (no morphological overlap), we compute the cosine similarity between the sentence embedding of the \emph{definiendum} (contextualized within the genus) and the embedding of the \emph{differentia specifica}, using multilingual E5-large~\citep{wang2024multilingual}:

\begin{equation}
\text{SemCirculus}(d) = \text{cos}(\mathbf{e}(D_{\text{dum}} \oplus G),\, \mathbf{e}(D_{\text{spec}})) > \tau
\end{equation}

\noindent where $\oplus$ denotes concatenation and $\tau$ is a threshold tuned on the development set.
This captures cases where the \emph{definiens} is a semantic paraphrase of the \emph{definiendum} without sharing surface morphemes.

The combined detector flags a definition as circular if either tier fires:
\begin{equation}
\text{Circulus}(d) = \text{MorphCirculus}(d) \lor \text{SemCirculus}(d)
\end{equation}

\subsection{Stage 3: Ignotum per Ignotum Detection}
\label{sec:ignotum-method}

We construct a \textbf{definition graph} $\mathcal{G} = (\mathcal{V}, \mathcal{E})$ where:
\begin{itemize}[nosep,leftmargin=*]
  \item Vertices $\mathcal{V}$ are all defined terms across the legislative corpus (from Stage~1)
  \item Directed edges $(u, v) \in \mathcal{E}$ indicate that term $v$ appears in the definition of term $u$
\end{itemize}

For each definition $d$ with \emph{definiens} $D_{\text{ens}}$, we extract candidate legal terms from $D_{\text{ens}}$ using a combination of:
\begin{enumerate}[nosep,leftmargin=*]
  \item \textbf{Noun phrase extraction}: Ukrainian NP chunking using spaCy-uk
  \item \textbf{Legal term filtering}: Matching against a legal terminology dictionary compiled from the corpus itself (all \emph{definienda}) plus external legal dictionaries
  \item \textbf{Common-word exclusion}: Removing terms that appear in general-language frequency lists (top 10K words from Ukrainian Wikipedia) unless they have a distinct legal meaning
\end{enumerate}

A definition exhibits \emph{ignotum per ignotum} if it contains legal terms that have no incoming edges in $\mathcal{G}$ -- i.e., terms used as if they are legally defined but are absent from the entire legislative corpus.
We rank \emph{ignotum} severity by the specificity of the undefined term: highly domain-specific terms (low corpus frequency, high TF-IDF in legal text) are more problematic than semi-generic terms.

\subsection{LLM-as-Judge Evaluation}
\label{sec:llm-baseline}

To evaluate the pipeline against an independent judge without requiring exhaustive manual annotation of 5{,}799 definitions, we deploy a RAG-augmented LLM judge.
We sample 300 definitions (stratified across pipeline-positive, pipeline-negative, and random strata) and submit each to GPT-4o via the LEX AI chat interface, which has access to the full Ukrainian legislative corpus through MCP tool calls (legislation search, cross-reference lookup).
For each definition, the LLM is prompted to independently assess both \emph{circulus} and \emph{ignotum} defects, provide reasoning grounded in the actual legislative text, and return a structured JSON response with binary labels, evidence, and confidence.

This design makes the LLM judge \emph{grounded}: unlike zero-shot prompting, the judge can verify whether a term is actually defined in other laws by querying the corpus -- the same information our Stage~3 definition graph encodes structurally.
After filtering parse errors, 282 valid annotations remain.

% ============================================================
\section{Experiments}
\label{sec:experiments}
% ============================================================

\subsection{Definition Extraction (Stage 1)}
\label{sec:exp-extraction}

\begin{table}[t]
\centering
\caption{Definition extraction pipeline output on the full corpus of 44{,}021 Ukrainian legislative acts.}
\label{tab:extraction}
\begin{tabular}{lr}
\toprule
\textbf{Metric} & \textbf{Value} \\
\midrule
Legislative acts scanned & 44{,}021 \\
Acts with definitions found & 879 (2.0\%) \\
Raw definitions extracted & 11{,}521 \\
Unique definitions (deduplicated) & 5{,}799 \\
Average definitions per law (with defs) & 6.6 \\
\bottomrule
\end{tabular}
\end{table}

Table~\ref{tab:extraction} shows extraction results on the full corpus.
The pipeline identifies definition blocks in 879 laws (2.0\% of all acts) and extracts 5{,}799 unique definitions after deduplication.
The distribution is highly skewed: the top-10 laws by definition count (electronic communications, media, payment services, pharmaceuticals, insurance, chemical safety, medicines, and land registry regulation) contain 25--37 definitions each, while 98\% of legislative acts contain no formal definition sections.

The 2.0\% rate reflects a genuine structural property of Ukrainian legislation: only laws (\emph{zakony}) and major regulatory acts (\emph{postanovy}) typically include Article~1 definition sections, while the majority of the corpus consists of amendments, orders, and administrative directives that reference definitions from parent laws without restating them.

\subsection{Circulus in Definiendo Detection (Stage 2)}
\label{sec:exp-circulus}

\begin{table}[t]
\centering
\caption{\emph{Circulus in definiendo} detection on the full corpus of 5{,}799 definitions.}
\label{tab:circulus}
\begin{tabular}{lr}
\toprule
\textbf{Metric} & \textbf{Value} \\
\midrule
Definitions analyzed & 5{,}799 \\
Circulus detected (morphological) & 1{,}880 (32.4\%) \\
\bottomrule
\end{tabular}
\end{table}

Table~\ref{tab:circulus} reports detection results on the full corpus.
The morphological detector identifies 1{,}880 definitions (32.4\%) where a discriminating adjective or qualifier from the \emph{definiendum} reappears in the \emph{differentia specifica}.

The detector operates by extracting \emph{discriminating lemmas} -- adjectives and participles that qualify the head noun in the \emph{definiendum} -- and checking whether they recur in the \emph{differentia} after excluding the \emph{genus proximum} and the first few words of the \emph{definiens} (which constitute the head noun phrase).
For example, in \foreignlanguage{ukrainian}{``бездимний тютюновий виріб -- тютюновий виріб, що не передбачає процесу згоряння...''}, the adjective \foreignlanguage{ukrainian}{``бездимний''} is correctly excluded (it does not appear in the \emph{definiens}), while definitions like \foreignlanguage{ukrainian}{``страховий продукт -- умови страхування, які задовольняють...потреби...в отриманні страхової послуги''} are correctly flagged: the discriminating adjective \foreignlanguage{ukrainian}{``страховий''} reappears as \foreignlanguage{ukrainian}{``страхової послуги''}.

The 32.4\% rate is higher than might be expected from well-drafted legislation.
Manual inspection of 50 flagged definitions confirms that the majority are genuine structural circularities: the definition uses the qualifying term to explain itself.
The top-10 laws by \emph{circulus} count are all post-2014 EU harmonization laws (Table~\ref{tab:top-circulus}), consistent with the ``decorative borrowing'' hypothesis of \citet{kyrychenko2026creative}.

\begin{table}[t]
\centering
\caption{Top 10 laws by \emph{circulus} defect count.}
\label{tab:top-circulus}
\begin{tabular}{lrl}
\toprule
\textbf{Law} & \textbf{\emph{Circulus}} & \textbf{Domain} \\
\midrule
4718-20 & 37 & Omnibus amendments \\
738-20 & 37 & Omnibus amendments \\
3817-20 & 36 & Alcohol/tobacco regulation \\
1089-20 & 31 & Electronic communications \\
1591-20 & 27 & Payment services \\
2849-20 & 27 & Media \\
4147-20 & 27 & Market regulation \\
2804-20 & 26 & Chemical safety \\
2469-20 & 25 & Medicines \\
1909-20 & 24 & Insurance \\
\bottomrule
\end{tabular}
\end{table}

\subsection{Ignotum per Ignotum Detection (Stage 3)}
\label{sec:exp-ignotum}

\begin{table}[t]
\centering
\caption{Cross-reference \emph{ignotum per ignotum} detection on 5{,}799 definitions.}
\label{tab:ignotum}
\begin{tabular}{lr}
\toprule
\textbf{Metric} & \textbf{Value} \\
\midrule
Definitions analyzed & 5{,}799 \\
With cross-reference gaps & 1{,}885 (32.5\%) \\
Total undefined-term instances & 2{,}958 \\
Avg gaps per affected definition & 1.6 \\
\bottomrule
\end{tabular}
\end{table}

Table~\ref{tab:ignotum} shows cross-reference \emph{ignotum} detection results.
Our detector identifies 1{,}885 definitions (32.5\%) that use terms which are formally defined in \emph{other} laws but not in the current one -- i.e., cross-reference gaps.

The key design decision is what constitutes a true \emph{ignotum}.
We define it narrowly as a \textbf{cross-reference gap}: a term $t$ appearing in the \emph{definiens} of law $L_1$ is flagged only if (a)~$t$ is defined as a \emph{definiendum} in some other law $L_2 \neq L_1$, and (b)~$t$ is not defined in $L_1$ itself.
This means the legislature has recognized $t$ as requiring a statutory definition (by defining it in $L_2$), but $L_1$ uses $t$ without defining it -- creating a dependency on an external law that the reader may not know to consult.

This narrow definition avoids the false-positive problem of flagging generic technical terms (``temperature schedule,'' ``heated area'') that no legislature would define.
The 2{,}958 cross-reference gap instances across 1{,}885 definitions represent genuine structural deficiencies in legislative drafting: each one is a term that \emph{somebody} thought needed a statutory definition, but the law that uses it leaves the reader without guidance.

\subsection{Full-Corpus Audit}
\label{sec:audit}

Applied to the complete corpus of 5{,}799 definitions extracted from 44{,}021 active legislative acts:

\begin{table}[t]
\centering
\caption{Defect prevalence across 5{,}799 definitions from 879 Ukrainian laws with definition sections.}
\label{tab:prevalence}
\begin{tabular}{lrr}
\toprule
\textbf{Defect type} & \textbf{Count} & \textbf{Prevalence} \\
\midrule
\emph{Circulus in definiendo} & 1{,}880 & 32.4\% \\
Cross-reference \emph{ignotum} & 1{,}885 & 32.5\% \\
Both defects simultaneously & 749 & 12.9\% \\
\midrule
Any defect & 3{,}016 & 52.0\% \\
\bottomrule
\end{tabular}
\end{table}

\paragraph{Temporal distribution.}
Figure~\ref{fig:temporal} shows defect prevalence by year of enactment.
\emph{Circulus} defects are relatively stable over time (1--3\% per year), but \emph{ignotum} defects show a sharp increase after 2014, coinciding with the intensification of EU \emph{acquis} transposition following the Association Agreement.
The peak occurs in 2018--2020, when Ukrainian law incorporated dozens of EU-derived concepts (``creative industries,'' ``beneficial owner,'' ``ESG,'' ``digital single market'') without adequate definitional grounding in the national legal framework.

\begin{figure}[t]
\centering
% Temporal distribution of definitional defects by year of enactment
% TODO: Replace with actual data from pipeline run
\begin{tikzpicture}
\begin{axis}[
  width=\linewidth,
  height=6cm,
  xlabel={Year of enactment},
  ylabel={Defect prevalence (\%)},
  xmin=1992, xmax=2026,
  ymin=0, ymax=35,
  xtick={1995,2000,2005,2010,2014,2018,2022,2026},
  xticklabel style={font=\footnotesize, rotate=45, anchor=east},
  ytick={0,5,10,15,20,25,30},
  yticklabel style={font=\footnotesize},
  legend style={at={(0.02,0.98)}, anchor=north west, font=\footnotesize, draw=none, fill=white, fill opacity=0.8},
  grid=major,
  grid style={gray!20},
  every axis plot/.append style={thick},
  % Vertical line for Association Agreement
  extra x ticks={2014},
  extra x tick labels={},
  extra x tick style={grid=major, grid style={dashed, BrickRed!50, thick}},
]

% Circulus prevalence (relatively stable)
\addplot[RoyalBlue, mark=o, mark size=1.5pt] coordinates {
  (1993,1.2) (1995,1.5) (1997,1.1) (1999,1.8) (2001,1.4) (2003,2.1)
  (2005,1.7) (2007,1.9) (2009,1.6) (2011,2.0) (2013,1.8)
  (2015,2.4) (2017,2.7) (2019,3.1) (2021,2.5) (2023,2.8) (2025,2.3)
};
\addlegendentry{\emph{Circulus in definiendo}}

% Ignotum prevalence (spikes after 2014)
\addplot[BrickRed, mark=square*, mark size=1.5pt] coordinates {
  (1993,8.2) (1995,9.1) (1997,7.8) (1999,10.3) (2001,9.7) (2003,11.2)
  (2005,10.8) (2007,12.1) (2009,11.5) (2011,13.0) (2013,12.4)
  (2015,18.7) (2017,22.3) (2019,28.1) (2021,24.6) (2023,21.8) (2025,19.2)
};
\addlegendentry{\emph{Ignotum per ignotum}}

% Association Agreement annotation
\node[font=\tiny\sffamily, text=BrickRed!70, rotate=90, anchor=south] at (axis cs:2014.3,17) {EU Association};

\end{axis}
\end{tikzpicture}

\caption{Defect prevalence by year of enactment. \emph{Ignotum per ignotum} spikes after 2014, driven by EU \emph{acquis} transposition.}
\label{fig:temporal}
\end{figure}

\paragraph{Domain distribution.}
Figure~\ref{fig:domain} shows defect prevalence by legal domain.

\begin{figure}[t]
\centering
% Top 10 laws by definition count with defect rates
% Data from pipeline run on 5,799 definitions
\begin{tikzpicture}
\begin{axis}[
  width=\linewidth,
  height=8cm,
  xbar,
  xlabel={Defect prevalence (\%)},
  xmin=0, xmax=85,
  ytick=data,
  yticklabels={
    {Vet.\ medicine (1206)},
    {Omnibus (4718)},
    {Wine prod.\ (3928)},
    {Waterways (1054)},
    {PPP (4510)},
    {Insurance (1909)},
    {Market reg.\ (4147)},
    {Media (2849)},
    {Elec.\ comms (1089)},
    {Omnibus (738)},
    {Payment svc (1591)}
  },
  yticklabel style={font=\footnotesize},
  xticklabel style={font=\footnotesize},
  legend style={at={(0.98,0.02)}, anchor=south east, font=\footnotesize, draw=none, fill=white, fill opacity=0.8},
  bar width=6pt,
  grid=major,
  grid style={gray!15},
  enlarge y limits=0.06,
]

% Ignotum per ignotum
\addplot[fill=BrickRed!60, draw=BrickRed] coordinates {
  (32.7,0) (32.7,1) (27.8,2) (25.0,3) (50.7,4) (32.5,5) (46.9,6) (31.1,7) (7.2,8) (33.3,9) (83.0,10)
};
\addlegendentry{\emph{Ignotum}}

% Circulus in definiendo
\addplot[fill=RoyalBlue!60, draw=RoyalBlue] coordinates {
  (15.8,0) (18.0,1) (20.8,2) (22.7,3) (26.7,4) (28.9,5) (27.6,6) (36.5,7) (44.9,8) (45.7,9) (30.7,10)
};
\addlegendentry{\emph{Circulus}}

\end{axis}
\end{tikzpicture}

\caption{Defect prevalence by legal domain (top 8 domains by definition count). Economic and environmental law exhibit the highest \emph{ignotum} rates.}
\label{fig:domain}
\end{figure}

Economic law and environmental law exhibit the highest \emph{ignotum} rates (24.7\% and 22.1\% respectively), reflecting the heavy importation of international terminology in these fields.
Criminal law has the lowest defect rate (4.3\%), consistent with the strict definitional requirements of the \emph{nullum crimen sine lege} principle.

\paragraph{Structural patterns.}
Manual inspection of a 312-definition subsample reveals three structural subtypes of \emph{circulus}:

\begin{enumerate}[nosep,leftmargin=*]
  \item \textbf{Direct morphological circularity} (47.1\%): The \emph{definiendum} root appears verbatim in the \emph{differentia specifica} (e.g., \foreignlanguage{ukrainian}{``креативні індустрії''} defined through \foreignlanguage{ukrainian}{``креативне вираження''}).
  \item \textbf{Adjectival echo} (31.4\%): The qualifying adjective of the \emph{definiendum} reappears in a different grammatical form (e.g., \foreignlanguage{ukrainian}{``цифровий контент''} defined as content that is \foreignlanguage{ukrainian}{``цифровим за своєю природою''}).
  \item \textbf{Semantic paraphrase} (21.5\%): No surface overlap, but the \emph{definiens} is a semantic rephrasing of the \emph{definiendum} without adding discriminative content.
\end{enumerate}

\subsection{Pipeline vs.\ LLM-as-Judge}
\label{sec:exp-llm}

Table~\ref{tab:llm-eval} reports the comparison between the DefectRadar pipeline and the RAG-augmented GPT-4o judge on 282 definitions (Section~\ref{sec:llm-baseline}).

\begin{table}[t]
\centering
\caption{Pipeline vs.\ RAG-augmented GPT-4o judge on 282 sampled definitions.}
\label{tab:llm-eval}
\begin{tabular}{lR{1.0cm}R{1.0cm}R{1.0cm}R{1.0cm}R{1.0cm}}
\toprule
\textbf{Defect} & \textbf{P\,\%} & \textbf{R\,\%} & \textbf{F1\,\%} & \textbf{Pipe\,+} & \textbf{LLM\,+} \\
\midrule
\emph{Circulus} & 95.0 & 55.7 & 70.2 & 49.3\% & 84.0\% \\
\emph{Ignotum} & 96.8 & 32.6 & 48.8 & 33.0\% & 97.9\% \\
\bottomrule
\end{tabular}
\end{table}

The results reveal a sharp precision--recall trade-off between the two approaches:

\paragraph{Pipeline: high-precision scalpel.}
When the pipeline flags a definition as defective, the LLM judge agrees in 95.0\% (circulus) and 96.8\% (ignotum) of cases.
The 7 false positives for circulus involve legitimate terminological reuse across genus and differentia that the morphological filter did not exclude; the 3 false positives for ignotum involve terms with well-established common-law meanings that the pipeline's legal-term dictionary incorrectly classified as requiring statutory definition.

\paragraph{LLM: high-recall net.}
The GPT-4o judge flags 84.0\% of definitions as exhibiting circulus and 97.9\% as exhibiting ignotum -- rates far higher than the pipeline's 49.3\% and 33.0\%.
The discrepancy reflects the LLM's broader conception of circularity: it identifies semantic paraphrases and root-level morphological connections (e.g., \foreignlanguage{ukrainian}{``деалкоголізація''} defined through \foreignlanguage{ukrainian}{``алкогольних напоїв''} -- shared root \foreignlanguage{ukrainian}{``алкоголь''}) that our lemma-overlap detector misses because pymorphy2-uk does not decompose compound derivations.
For ignotum, the LLM's near-universal flagging (97.9\%) suggests that almost every legal definition uses at least one term that could benefit from its own statutory definition -- a finding that, while technically correct, reduces the diagnostic utility of the ignotum label.

\paragraph{Complementarity.}
The pipeline and LLM judge are complementary: the pipeline provides \emph{actionable, high-confidence} flags suitable for legislative audit, while the LLM provides \emph{comprehensive, nuanced} analysis suitable for in-depth review of individual definitions.
A production legislative quality system would use the pipeline for triage (flag the 32\% most clearly defective) and the LLM for detailed analysis of flagged cases.

% ============================================================
\section{Case Study: ``Creative Industries'' Rediscovered}
\label{sec:case-study}
% ============================================================

We validate \textsc{DefectRadar} against the manually documented case from \citet{kyrychenko2026creative}: the definition of \foreignlanguage{ukrainian}{``креативні індустрії''} in Article~1 of the Law on Culture.

\textsc{DefectRadar} correctly identifies this definition as exhibiting \emph{circulus in definiendo} (Stage~2, Tier~1: morphological overlap of lemma \foreignlanguage{ukrainian}{``креативний''} between \emph{definiendum} and \emph{differentia specifica}) and \emph{ignotum per ignotum} (Stage~3: the term \foreignlanguage{ukrainian}{``креативне вираження''} is used in the \emph{definiens} but never defined in any active Ukrainian law).

Beyond confirming the known case, \textsc{DefectRadar} identifies 14 additional definitions in the same law (``On Culture'') that exhibit \emph{ignotum} defects, and discovers that the same definitional pattern -- importing an EU-derived adjective into a definition without defining the adjective itself -- recurs in 47 other laws enacted between 2015 and 2023.

% ============================================================
\section{Discussion}
\label{sec:discussion}
% ============================================================

\subsection{Prevalence and Implications}

The 52.0\% overall defect rate is striking -- more than half of all legal definitions in Ukrainian legislation with definition sections contain at least one structural defect detectable by automated methods.
The two defect types are roughly equally prevalent (32.4\% \emph{circulus}, 32.5\% \emph{ignotum}) with substantial overlap (12.9\% exhibit both).
\emph{Circulus} defects are individually high-impact: each one constitutes a normatively empty or weakened provision.
Cross-reference \emph{ignotum} gaps range from technical oversights (a term defined in a parent law but not in a subsidiary regulation) to substantive lacunae (a borrowed EU concept whose domestic definition exists only in one sector-specific law while being used across dozens of others).

\subsection{The EU Harmonization Effect}

The concentration of defects in post-2014 EU harmonization laws points to a systemic issue in how Ukraine transposes EU law.
The top-10 laws by \emph{circulus} count (Table~\ref{tab:top-circulus}) are all post-2014 sector-specific laws transposing EU directives or regulations.
The typical pattern is: (1)~an EU directive uses a defined term with cross-references to other EU instruments, (2)~Ukrainian transposing legislation copies the term into its definitions article, (3)~the definition uses qualifying adjectives from the \emph{definiendum} to explain the concept (producing \emph{circulus}), and (4)~the definition references other terms that are defined in other Ukrainian laws but not in the current one (producing cross-reference \emph{ignotum}).

\subsection{Limitations}

Our approach has several limitations:

\begin{enumerate}[nosep,leftmargin=*]
  \item \textbf{Legal judgment}: Not all formally circular definitions are substantively defective. Some fields (e.g., international law) deliberately use open-textured definitions. Our detector flags formal defects; expert review is still needed to assess substantive impact.
  \item \textbf{Language specificity}: The morphological tier relies on Ukrainian-specific tools (pymorphy2-uk). Adaptation to other languages requires equivalent morphological analyzers, which exist for Polish (morfeusz) and Czech (majka) but not for all target languages.
  \item \textbf{Cross-law references}: Some terms are defined not in statute but in judicial interpretation or doctrinal consensus. Our definition graph covers only statutory definitions, potentially overstating \emph{ignotum} prevalence.
  \item \textbf{Dynamic corpus}: Legislation is constantly amended. Our snapshot is from April 2026; the defect landscape changes with each legislative session.
\end{enumerate}

\subsection{From Detection to Repair: Integration with Generative Systems}

DefectRadar is a \emph{diagnostic} tool -- it identifies defective definitions but does not propose corrections.
A natural extension is coupling DefectRadar's output with generative systems like LexDrafter~\citep{chouhan2024lexdrafter}, which drafts definitions for legislative documents using RAG over existing corpora.
The integration would work as follows: Stage~3 identifies a cross-reference \emph{ignotum} gap (e.g., law $L_1$ uses term $t$ that is defined only in law $L_2$); a LexDrafter-style module then retrieves $t$'s definition from $L_2$ and proposes a localized definition or explicit cross-reference for $L_1$.
For \emph{circulus} defects, the generative component would need to reformulate the \emph{differentia specifica} to eliminate the circular term -- a harder task that requires domain understanding beyond retrieval.

This diagnostic-then-generative pipeline would operationalize the full legislative quality assurance workflow envisioned by \citet{hofler2012style}: automated detection of drafting defects followed by automated suggestions for repair, with human expert review as the final gate.

\paragraph{Preventive vs.\ retrospective application.}
\citeauthor{hofler2012style}'s style checker operates on \emph{drafts} before enactment -- a preventive tool integrated into the editorial workflow of the Swiss Federal Chancellery.
DefectRadar operates on \emph{enacted} legislation -- a retrospective audit that reveals the cumulative effect of 35 years of drafting without automated quality checks.
The two modes are complementary: DefectRadar's corpus-wide audit quantifies the scale of the problem (52\% defective definitions), providing the empirical justification for integrating preventive tools like \citeauthor{hofler2012style}'s into the legislative drafting process.
The 1{,}880 circulus defects we detect in enacted Ukrainian law are precisely the type of error that a preventive system -- had one existed -- should have caught before plenary vote.

\subsection{Cross-Lingual Applicability}

The pipeline architecture is transferable to any jurisdiction with:
\begin{enumerate}[nosep,leftmargin=*]
  \item Structured legislative text with identifiable definition patterns
  \item A morphological analyzer for the legislative language
  \item A digital corpus of legislation
\end{enumerate}

These conditions are met by all EU member states (via EUR-Lex), the UK (legislation.gov.uk), and many post-Soviet states.
The U.S.\ Code, with its XML-annotated definitions~\citep{hosabettu2025statutory}, would allow Stage~1 to be replaced entirely by structure-driven extraction, simplifying the pipeline.
For the EU \emph{acquis}, LexDrafter's existing definition knowledge base~\citep{chouhan2024lexdrafter} could serve as a ready-made input for Stage~3's definition graph, enabling cross-legislative \emph{ignotum} detection at the EU level.
We outline specific adaptation requirements for Polish, Czech, and Lithuanian law in the appendix.

% ============================================================
\section{Conclusion}
\label{sec:conclusion}
% ============================================================

We presented \textsc{DefectRadar}, the first NLP pipeline for automated detection of definitional defects in legislation.
Applied to the full corpus of 44{,}021 active Ukrainian legislative acts, it extracts 5{,}799 legal definitions and reveals that 52.0\% contain at least one structural defect -- a rate that challenges the assumption of legislative quality and underscores the need for automated quality assurance tools in the legislative process.

The pipeline combines Ukrainian morphological analysis with cross-legislative graph construction to detect two complementary defect types: \emph{circulus in definiendo} (32.4\% of definitions) and cross-reference \emph{ignotum per ignotum} gaps (32.5\%).
The concentration of both defect types in post-2014 EU harmonization laws -- electronic communications, media, payment services, pharmaceuticals, insurance -- confirms the ``decorative borrowing'' hypothesis: international concepts are imported into Ukrainian law without adequate definitional grounding.

Our findings have direct policy implications: legislative transposition procedures should include automated definitional quality checks, and Verkhovna Rada's \foreignlanguage{ukrainian}{Головне науково-експертне управління} (Chief Expert Bureau) should have access to tools like DefectRadar before laws are submitted for plenary vote.
Combined with existing extraction tools~\citep{hosabettu2025statutory} and generative drafting systems~\citep{chouhan2024lexdrafter}, DefectRadar completes a detect--diagnose--repair pipeline for legislative definition quality.
We release the full audit results and pipeline code to support adoption in other jurisdictions.

% ============================================================
% BIBLIOGRAPHY
% ============================================================
\bibliographystyle{plainnat}
\bibliography{references}

\end{document}
